% SPDX-License-Identifier: CC-BY-SA-4.0
% Author: Matthieu Perrin
% Part: 
% Section: 
% Sub-section: 
% Frame: 

\begingroup

\begin{frame}[fragile]{Démonstration}

  \vspace{-22mm}
  \begin{lstlisting}[numbers=left, gobble=4]
    public void push(T value) {
      do {
        Node old = head.get();
      } while(!head.compareAndSet(old, new Node(value, old)));
    }

    public T pop() {
      for(;;) {
        Node old = head.get();
        if(old==null) return null;
        if(head.compareAndSet(old, old.next())
        return old.value();
      }
    }
  \end{lstlisting}

  \obBlock<1>[anchor=north, y=-5mm]{Preuve de vivacité}{
    \begin{itemize}
    \item L'algorithme n'est pas wait-free :
      \begin{itemize}
      \item $p_1$ et $p_2$ appellent \lstinline{push} continuellement ;
      \item $p_2$ peut perdre tous ses  \lstinline{compareAndSet} ligne 4.
      \end{itemize}
    \item L'algorithme est lock-free
      \begin{itemize}
      \item Si $p_i$ ne termine pas lors d'un tour de boucle (ligne 4 ou 11),
      \item Alors \lstinline{head} a changé depuis le \lstinline{get} (ligne 3 ou 9);
      \item Donc un $p_j$ a écrit dans \lstinline{head} (ligne 4 ou 11) et terminé une opération.
      \end{itemize}
    \end{itemize}
  }

  \obBlock<2>[anchor=north, y=-5mm]{Preuve de linéarisabilité}{
    \begin{itemize}
    \item Point de linéarisation de \lstinline{push} :
      \begin{itemize}
      \item dernier appel à \lstinline{compareAndSet} ligne 4
      \end{itemize}
    \item Point de linéarisation de \lstinline{pop} :
      \begin{itemize}
      \item dernier appel à \lstinline{get} ligne 9 \structure{si la pile est vide} 
      \item dernier appel à \lstinline{compareAndSet} ligne 11 \structure{sinon}
      \end{itemize}
    \item Entre les points de linéarisation, même état qu'en séquentiel
    \end{itemize}
  }

\end{frame}

\endgroup
\endinput
